\section{Conditioning on length from logarithmic finite estimates}
\label{R:cal:section}

We now obtain exact-length bridge laws from weaker finite length
information: a logarithmic window and a first moment restricted to bridges
whose width is at most a fixed multiple of their height. Fourier estimates first produce a whole interval
of renewal masses with a common lower bound. Fresh independent blocks then
move the remaining target length into that interval at subpolynomial cost.
This route also gives the neighboring-length ratio and local convergence
of the finite bridge laws. The later geometric estimates and unrestricted
changes of law use the finite inputs stated here, but do not enter this
conditioning argument.

Throughout this section the critical weight is $\rho^L$, where $L$ counts
visited centers and boundary ports carry no weight. The strict bridge starts
at one fixed port and its terminal port is free. Height $H$ is a positive
integer; one band has Euclidean height $\sqrt3/2$. Its transverse displacement
is denoted by $X$, with $2X$ integral on the row-staggered lattice. Write $D$
for its diameter. The irreducible law $p$ and its independent concatenation
are those of Proposition~\ref{R:renew:factorization}. Dependence between $H,X,L,D$
within one irreducible is allowed. Put
\[
 a=\frac34,\qquad \vartheta=1-a=\frac14,\qquad
 d_* =\frac43,\qquad \alpha=\frac{a}{d_*}=\frac9{16}.
\]
We use $\vartheta$ for the boundary exponent to keep it distinct from the
height-tail exponent $a$.

\subsection{The finite estimates and a logarithmic renewal denominator}
\label{R:cal:conditioning}
The finite input is supplied by the calibrated cylinder analysis
\cite[Theorems~1.1--1.2 and Proposition~11.2]{compL}.
We state its precise form here, so that the conditioning argument can be
read without revisiting that representation.

\begin{proposition}[Finite calibrated input]\label{R:cal:inputs}
Let $B_h$ be the critical mass of strict bridges of height $h$ from one
fixed port, and let $B_0=1$. Then $B_h\asymp(1+h)^{-1/4}$. For sufficiently
large fixed $T$, bridges confined to transverse width $Th$ have mass at
least $c_T h^{-1/4}$ and first-length mass at most $C_T h^{13/12}$.
The arch-plus-bridge mass in a height-$h$ strip with diameter larger than
$r\ge Ch$ is at most $C h^{-1/4}e^{-cr/h}$. Since a self-avoiding path of
diameter $r$ has at most $Cr^2$ centers, this estimate also controls every
fixed length moment in the tail. In particular the mass at
$D>h\log^2 h$, including any fixed length power, is smaller than every
inverse power of $h$.

The half-plane arch mass at prescribed horizontal gap $g$ satisfies
$A_g\asymp g^{-5/4}$; its lower bound remains valid with diameter at most
$Cg$ for a sufficiently large fixed $C$.
The normalized irreducible law satisfies
\[
 1-\mathbb E_p e^{-tH}\asymp t^{3/4}\quad(0<t\le1),\qquad
 p(D>r)\le C r^{-3/4}\quad(r\ge1).
\]
For $h$ large, put $R_0=h(\log h)^{1/64}$ and
\[
 l_-=h^{4/3}(\log h)^{-1/8},\qquad
 l_+=h^{4/3}(\log h)^{1/2}.
\]
For a fixed finite $C$,
\begin{equation}
 \sum_{j\le2R_0}B_j\{l_-\le L\le2l_+\}
       \ge h^{3/4}(\log h)^{-C}.
 \label{R:cal:window}
\end{equation}
The estimate sums over heights and terminal ports; it does not prescribe
either one of them.
\end{proposition}

\begin{lemma}[Length deficit]\label{R:cal:deficit}
There is a finite $P$ such that, for all sufficiently small $t>0$,
\[
 c t^\alpha(\log(1/t))^{-P}
 \le q_L(t):=1-\mathbb E_p e^{-tL}\le C t^\alpha.
\]
\end{lemma}
\begin{proof}
Take $t=h^{-4/3}$. On the length window in \eqref{R:cal:window},
$1-e^{-tL}\ge c(\log h)^{-1/8}$. Summing this damping loss over the
$O(h(\log h)^{1/64})$ allowed heights shows that some height
$1\le h_*\le2R_0$ has loss at least $h^{-1/4}(\log h)^{-C'}$.
For a bridge with irreducible lengths $L_i$,
$1-e^{-t\sum_iL_i}\le\sum_i(1-e^{-tL_i})$.
The marked reward identity of Lemma~\ref{R:renew:reward} therefore bounds
this loss above by
\[
 \sum_{j\le h_*}\mathbb E_p[1-e^{-tL};H=j]
                 \sum_{b+d=h_*-j}B_bB_d
 \le C h_*^{1/2}q_L(t).
\]
Here $B_j\asymp(1+j)^{-1/4}$ gives the convolution bound. Since
$h_*\le2h(\log h)^{1/64}$, we obtain
\[
 q_L(h^{-4/3})\ge c h^{-3/4}(\log h)^{-P}
\]
for some finite $P$, proving the lower estimate.

For the upper estimate fix a sufficiently large confinement constant $T$
in Proposition~\ref{R:cal:inputs}. At each height $d\le h$, its confined
bridges have mass at least $c_Td^{-1/4}$ and first-length mass at most
$C_Td^{13/12}$. Choosing a large fixed $C'$ and restricting to
$L\le C'h^{4/3}$ preserves at least half the mass: the discarded fraction
is at most a constant times $d^{4/3}/(C'h^{4/3})$. On the retained paths
$e^{-tL}\ge e^{-C'}$. Summing over $d\le h$ gives damped bridge mass at
least $c h^{3/4}$. By renewal the total damped bridge mass over all heights
is $1/q_L(t)$, so $q_L(h^{-4/3})\le Ch^{-3/4}$.
\end{proof}

The upper deficit also gives the length tail and truncated mean:
\[
 p(L>x)\le\frac{q_L(1/x)}{1-e^{-1}}\le Cx^{-\alpha},\qquad
 \mathbb E_p[L;L\le x]\le Cx^{1-\alpha}\quad(x\ge1).
\]
Lemma~\ref{R:renew:growth} therefore gives both bounds for length and
height, and the upper bound for diameter. Writing
$S_Y(k)=\sum_{i=1}^kY_i$ for each coordinate,
we obtain almost surely
\[
 S_H(k)=k^{4/3+o(1)},\qquad S_L(k)=k^{16/9+o(1)},\qquad
 S_D(k)\le k^{4/3+o(1)}.
\]
Later slabs cannot re-enter below a completed height. The interpolation
between consecutive complete pieces in Theorem~\ref{R:renew:infinite}
therefore proves the same infinite-path conclusions from these weaker
inputs: mass $r^{4/3+o(1)}$ in a radius-$r$ ball and distance
$n^{3/4+o(1)}$ at center time $n$.

\begin{theorem}[Exact compatible lengths from a logarithmic window]
\label{R:cal:exact}
Let $u_n=p(\exists k:S_L(k)=n)$. For every sufficiently large even $n$,
\[
 u_n\ge n^{-7/16-o(1)},\qquad \frac{u_{n+2}}{u_n}\longrightarrow1.
\]
Conditioning the independent sequence on hitting $n$ gives the uniform
strict bridge of length $n$, with height and terminal port free. Under this
law the height, endpoint distance, and maximum radius are
$n^{3/4+o(1)}$ in probability. The finite bridge laws converge locally to
the independent-irreducible law.
\end{theorem}
\begin{proof}
\emph{A block law from the finite window.}
The height deficit gives
\[
 p\{S_H(k)\le2R_0\}\le e\exp(-c kR_0^{-a}).
\]
Consequently, the contribution of counts $k>h^a\log^2h$ to
\eqref{R:cal:window} is negligible. Dividing the remaining window mass by
this number of counts yields a block whose probability of lying in the
length window is bounded below by a negative power of $\log h$.
Since $l_-/l_+=(\log h)^{-5/8}$, fix any $b\in(5/8,1)$, say $b=3/4$.
For each sufficiently large $r$, choose $h$ so that $l_+$ is comparable
to $r$ and $2l_+\le r/2$. Then there is a deterministic integer $k(r)$
such that
\begin{equation}
 p\{r(\log r)^{-b}\le S_L(k(r))\le r/2\}
       \ge(\log r)^{-C}.
 \label{R:cal:landing}
\end{equation}
Fix one such choice for every $r$. We will use these blocks to advance
through a large remaining gap without passing its target.

\paragraph{A whole interval of local renewal mass.}
Port lengths are even, and two- and four-center one-row irreducibles show
that $L/2$ has span one. Let $\phi$ be its characteristic function. The
length deficit with logarithmic loss implies, near zero,
\[
 1-|\phi(\theta)|^2
       \ge c|\theta|^\alpha(\log(1/|\theta|))^{-P'}.
\]
To verify this estimate, set $T=1/|\theta|$ and
$z=T/(\log T)^K$. The deficit bounds give
$\mathbb E_p\min(1,L/z)\ge c z^{-\alpha}(\log T)^{-P}$.
Lengths above $T$ contribute at most $CT^{-\alpha}$; those below
$z/(\log T)^{K'}$ contribute at most
$Cz^{-\alpha}(\log T)^{-K'(1-\alpha)}$, by the truncated mean bound.
Choose $K$ and then $K'$ large enough that these contributions leave a
fixed fraction of the lower bound. On the remaining range the phase
$\theta(L-2)/2$ has magnitude between a negative power of $\log T$ and
a fixed constant less than one. In
$1-|\phi(\theta)|^2=\mathbb E_p[1-\cos(\theta(L-L')/2)]$,
keep this range and the atom $L'=2$. The cosine bound proves the claim.
Away from zero, the atoms at $L/2=1,2$ give $|\phi|<1$.

Fourier inversion now gives, uniformly in the target and for every $k\ge1$,
\begin{align}
 \sup_j p(S_L(k)=j)&\le C(1+k)^{-1/\alpha}\log^C(2+k),
 \label{R:cal:length-atom}\\
 \sup_{j\text{ even}}|p(S_L(k)=j+2)-p(S_L(k)=j)|
 &\le C(1+k)^{-2/\alpha}\log^C(2+k).
 \label{R:cal:length-difference}
\end{align}
Indeed split the Fourier integral at
$k^{-1/\alpha}(\log k)^M$ for large $k$. The inner interval has that
length; for the difference its integrand has the additional factor
$|e^{i\theta}-1|\le|\theta|$. For sufficiently large fixed $M$, the
remaining integral near zero is smaller than any prescribed power of
$k^{-1}$, and the part away from zero is exponentially small. Increasing
the constants handles bounded $k$.

For a small fixed $c_0>0$, put
\[
 f_n(j)=\sum_{c_0n^\alpha\le k\le2c_0n^\alpha}p(S_L(k)=j).
\]
The upper deficit bound gives
\[
 p\{S_L(k)>n/3\}
 \le\frac{1-(1-q_L(1/n))^k}{1-e^{-1/3}}
 \le Ck n^{-\alpha}\le\tfrac12
\]
when $c_0$ is sufficiently small. Thus
$\sum_{0\le j\le n/3}f_n(j)\ge c n^\alpha$, and some even
$j_0\le n/3$ satisfies $f_n(j_0)\ge c n^{\alpha-1}$.
By \eqref{R:cal:length-difference},
\[
 |f_n(j+2)-f_n(j)|\le Cn^{\alpha-2}(\log n)^{C_1}.
\]
This preserves half the lower bound on an even interval
$[m,m+w]=[j_0,j_0+w]\subset[0,n/2]$, where
$w\ge c'n(\log n)^{-C_1}$. Positive increments make the events
$S_L(k)=j$ disjoint in $k$ on one independent sequence. Hence
$u_j\ge f_n(j)\ge c n^{\alpha-1}$ throughout this interval.

\paragraph{Landing before the final renewal hit.}
Aim for a partial sum in $[n-m-w,n-m]$. At a current sum
$x<n-m-w$, put $r=n-m-x$ and expose the next $k(r)$ fresh pieces.
On the event in \eqref{R:cal:landing}, the new remaining gap is between
$r/2$ and $r-r(\log r)^{-b}$. In particular the target $n-m$ is not
passed. Stop as soon as the gap is at most $w$, so the partial sum lies
in the required landing interval.

Before stopping every gap lies in $[w,n]$, and its logarithm is comparable
to $\log n$. At most
$J=O((\log n)^b(1+\log\log n))$ successful stages suffice, since each
reduces the gap by a fraction at least $(\log n)^{-b}$ and
$n/w\le C(\log n)^{C_1}$. The choice $k(r)$ depends only on the exposed
past, so every block end is a stopping time and the unused pieces remain
independent with law $p$. At each unfinished stage the conditional success
probability is at least $(\log n)^{-C}$. Declaring all later stages
successful once landing has occurred gives
\[
 p\{\text{all required stages succeed}\}
 \ge\exp\{-C(\log n)^b(1+\log\log n)^2\}=n^{-o(1)}.
\]
The fixed inequality $b<1$ is essential here. At the landing time the
remaining length to $n$ belongs to $[m,m+w]$. The fresh continuation
therefore hits $n$ with probability at least $c n^{\alpha-1}$, and
\[
 u_n\ge n^{\alpha-1-o(1)}.
\]

\paragraph{Conditioned count and spatial bounds.}
Every length-$n$ bridge has critical weight $\rho^n$, so conditioning the
sequence on hitting $n$ gives exactly the uniform strict-bridge law.
We first show that its number of pieces lies between
$n^{\alpha-\epsilon}$ and $n^{\alpha+\epsilon}$ with high probability,
for each fixed $0<\epsilon<\min(\alpha,1-\alpha)$; larger tolerances follow at once.

The tail and truncated mean give
$p\{S_L(j)>t\}\le Cj t^{-\alpha}$ by the union bound and Markov's
inequality. For $k\ge2$, split the pieces into two comparable groups.
If $S_L(k)=n$, one group has length at least $n/2$, at cost
$Ck n^{-\alpha}$; the independent other group pays at most
$Ck^{-1/\alpha}\log^C(2+k)$ for the residual length, by
\eqref{R:cal:length-atom}. Consequently
\[
 \sum_{2\le k\le n^{\alpha-\epsilon}}p(S_L(k)=n)
 \le Cn^{-\alpha}\sum_{k\le n^{\alpha-\epsilon}}
           k^{1-1/\alpha}\log^C(2+k)
 =n^{\alpha-1-\epsilon(2-1/\alpha)+o(1)}=o(u_n).
\]
Here $2-1/\alpha=2/9>0$. The one-piece term is at most
$Cn^{-\alpha}=o(u_n)$ because $2\alpha-1=1/8>0$.
For the upper cutoff, the lower deficit gives
\[
 p(S_L(k)=n)\le e\exp\{-c k n^{-\alpha}(\log n)^{-P}\}.
\]
Summing over $k>n^{\alpha+\epsilon}$ is negligible; there are at most $n$
possible counts since lengths are positive.

It suffices to take $0<\xi<3/4$. Choose $\epsilon>0$ sufficiently small. At the retained
counts, the height deficit gives
\[
 p\{S_H(k)<n^{3/4-\xi}\}
 \le\exp\{1-c n^{a\xi-\epsilon}\},
\]
which remains negligible after summing counts and dividing by $u_n$ if
$\epsilon<a\xi$. For the upper spatial bound put $r=n^{3/4+\xi}$.
The diameter tail similarly gives $p\{S_D(j)>r\}\le Cj r^{-a}$.
If $S_D(k)>r$, one of the two groups has diameter sum greater than $r/2$.
Test this event in that group and the exact residual length in the other.
By independence between groups and \eqref{R:cal:length-atom},
\[
 \begin{aligned}
 \sum_{n^{\alpha-\epsilon}\le k\le n^{\alpha+\epsilon}}
       p\{S_D(k)>r,\ S_L(k)=n\}
 &\le Cr^{-a}\sum_{k\le n^{\alpha+\epsilon}}
                  k^{1-1/\alpha}\log^C(2+k)\\
 &=n^{\alpha-1-a\xi+\epsilon(2-1/\alpha)+o(1)}=o(u_n)
 \end{aligned}
\]
for sufficiently small $\epsilon$. Neither split assumes independence
between the coordinates of one piece. The total height bounds endpoint
distance below, and $S_D(k)$ bounds the whole trace above. Absorbing fixed
geometric constants by slightly smaller tolerances proves the spatial
claims.

\paragraph{The neighboring-length ratio and local limit.}
In $u_{n+2}-u_n$, counts outside
$[n^{\alpha-\epsilon},n^{\alpha+\epsilon}]$ contribute $o(u_n)$ by the
same absolute bounds for both target lengths. On the retained counts,
\eqref{R:cal:length-difference} bounds the sum of successive differences by
\[
 C\sum_{k\ge n^{\alpha-\epsilon}}
       k^{-2/\alpha}\log^C(2+k)
 =n^{(\alpha-\epsilon)(1-2/\alpha)+o(1)}=o(u_n)
\]
for small $\epsilon$. Thus $u_{n+2}/u_n\to1$ on even lengths.
For a fixed tuple of $k$ irreducibles of total length $l$, its initial-tuple
probability under the length-$n$ bridge law is
$\rho^l u_{n-l}/u_n$ for $n>l$, which tends to $\rho^l$.
For fixed $k$ these limiting probabilities sum to one. A sufficiently
large $k$ determines any prescribed finite initial path, or the path
inside a fixed ball, since all later pieces lie above the completed
height. Finite-set truncation of the tuples proves the asserted local
convergence.
\end{proof}

\paragraph{Endpoint conventions and the half-plane limit.}
The half-plane paths considered here start at the inward center adjacent to a
fixed wall port and otherwise stay strictly above that wall; their terminal
center is free. They are precisely the paths in
Theorem~\ref{R:renew:uniform}. Its slot-flip ratio and the disjoint-prefix
factorization therefore prove that their uniform laws converge locally to
our independent-irreducible law. The proof applies in center length; a
bounded change to edge length does not change its conclusion.

The exact-length conclusion also transfers to uniform \(n\)-edge vertex bridges from a fixed vertex, with endpoints the height extrema (ties allowed, or take a strict initial minimum). Extend each to a strict midport bridge by these same bounded endpoint extensions, of length \(m=n+1,n+2\) or \(n+3\), with bounded multiplicity and start-port choices. At even \(m\) the exact-length estimates above give mass of bridges violating either spatial cutoff with fixed slack at most \(m^{\alpha-1-c(\xi)}\), \(c(\xi)>0,\ \alpha=9/16\) (both the diameter and small-piece-count bounds have strict power savings). Constant endpoint adjustments preserve this with slightly reduced slack. Conversely prepend to a strict bridge one down center below its start port, or for prescribed up initial type also one lower adjacent up center before this down center, as vertex root; append one outward up center through the top port or not. After translation this realizes \(n=m+\varepsilon_-+\varepsilon_+\) with \(\varepsilon_-\in\{0,1\}\) determined by initial type and \(\varepsilon_+\in\{0,1\}\) freely chosen so every large \(n\) uses even \(m\); multiplicity and weight losses are bounded. (If both extrema must be strict, append one or two outward centers instead.) The total vertex-bridge mass is thus \(\ge n^{\alpha-1-o(1)}\), proving the same typical exponent \(3/4\) there.
